\subsection{Inverses of Functions} 
We are especially interested in finding the inverse of a function, when that inverse is itself a function.
\begin{itemize}
\item The inverse of a function is always a \makeblank.
\item The inverse of a one-to-one relation is always a \makeblank.
\end{itemize}
To find the inverse of a function (given in function notation) we can go back to relation notation.
\begin{example}
Let $f(x)=3x-1$. Find $f\inv$.
\begin{lpic}{}{9}
\writethink{-1}{9}
\end{lpic}
\end{example}
\noindent In practice, we usually don't write relation notation, and use the following procedure.
\begin{procedure}{Finding an Inverse Function}
Use the following steps.
\begin{enumerate}
\item ``Swap $x$ and $y$.''
\item Solve for $y$.
\item Write in function notation.
\end{enumerate}
\end{procedure}